\documentclass[10pt,a4paper]{amsart}
\usepackage{amsmath,amssymb,amsthm,graphicx}
% Retained preamble for the Stacks Project chapter snapshot in this directory.
% The chapter's own preamble.tex is a copy of the shared Stacks Project preamble; the block
% below is sliced MECHANICALLY, line for line, from it. Nothing was added, removed or
% redefined; the wrapper class, the xr-hyper/hyperref wiring and the externaldocument list
% are omitted because the assembler supplies the document class.
%
% lines taken from preamble.tex: 3, 5, 9, 11, 12, 156, 157, 158, ... (51 lines total)
%
% package coverage restored from the chapter's own preamble (2026-10-13): xy, xy-2cell

\usepackage[all]{xy}
\xyoption{2cell}
\UseAllTwocells

\newenvironment{reference}{\comment}{\endcomment}
\newenvironment{slogan}{\comment}{\endcomment}
\newenvironment{history}{\comment}{\endcomment}
\theoremstyle{plain}
\newtheorem{theorem}[subsection]{Theorem}
\newtheorem{proposition}[subsection]{Proposition}
\newtheorem{lemma}[subsection]{Lemma}
\theoremstyle{definition}
\newtheorem{definition}[subsection]{Definition}
\newtheorem{example}[subsection]{Example}
\newtheorem{exercise}[subsection]{Exercise}
\newtheorem{situation}[subsection]{Situation}
\theoremstyle{remark}
\newtheorem{remark}[subsection]{Remark}
\newtheorem{remarks}[subsection]{Remarks}
\numberwithin{equation}{subsection}
\def\lim{\mathop{\mathrm{lim}}\nolimits}
\def\colim{\mathop{\mathrm{colim}}\nolimits}
\def\Spec{\mathop{\mathrm{Spec}}}
\def\Hom{\mathop{\mathrm{Hom}}\nolimits}
\def\Ext{\mathop{\mathrm{Ext}}\nolimits}
\def\SheafHom{\mathop{\mathcal{H}\!\mathit{om}}\nolimits}
\def\SheafExt{\mathop{\mathcal{E}\!\mathit{xt}}\nolimits}
\def\Sch{\mathit{Sch}}
\def\Mor{\mathop{\mathrm{Mor}}\nolimits}
\def\Ob{\mathop{\mathrm{Ob}}\nolimits}
\def\Sh{\mathop{\mathit{Sh}}\nolimits}
\def\NL{\mathop{N\!L}\nolimits}
\def\CH{\mathop{\mathrm{CH}}\nolimits}
\def\proetale{{pro\text{-}\acute{e}tale}}
\def\etale{{\acute{e}tale}}
\def\QCoh{\mathit{QCoh}}
\def\Ker{\mathop{\mathrm{Ker}}}
\def\Im{\mathop{\mathrm{Im}}}
\def\Coker{\mathop{\mathrm{Coker}}}
\def\Coim{\mathop{\mathrm{Coim}}}
\DeclareMathSymbol{\boxtimes}{\mathbin}{AMSa}{"02}
\def\QCohstack{\mathcal{QC}\!\mathit{oh}}
\def\Cohstack{\mathcal{C}\!\mathit{oh}}
\def\Spacesstack{\mathcal{S}\!\mathit{paces}}
\def\Quotfunctor{\mathrm{Quot}}
\def\Hilbfunctor{\mathrm{Hilb}}
\def\Curvesstack{\mathcal{C}\!\mathit{urves}}
\def\Polarizedstack{\mathcal{P}\!\mathit{olarized}}
\def\Complexesstack{\mathcal{C}\!\mathit{omplexes}}
\def\Pic{\mathop{\mathrm{Pic}}\nolimits}
\def\Picardstack{\mathcal{P}\!\mathit{ic}}
\def\Picardfunctor{\mathrm{Pic}}
\def\Deformationcategory{\mathcal{D}\!\mathit{ef}}
% lane-s3 patch-sd-envs.tex - DO NOT REMOVE LINES 16-21 (the executable patch).
% Contract (Chief-confirmed 2026-10-14):
%  * Lines 16-18: reference/slogan/history are declared as \comment containers in
%    every retained-preamble of this lane; the patch empties them so navigation
%    pointer text inside statements typesets without "Environment undefined".
%  * Line 21: \cite is globally no-op'd because the corpus rule is "cards carry
%    no bibliography": there is no bibliography backend at assembly time, and
%    every patched card's statement+proof span was scanned zero-cite before
%    patching (2026-10-14, lane-s3) - the only \cite occurrences live inside
%    the emptied reference environments (e.g. \cite{EGA} pointers). Removing
%    line 21 breaks compilation of every card whose statement carries a
%    reference environment with a citation (co01 wave: Citation `EGA' undefined).
%  * The patch is baked into task.json at card-build time (preamble_span
%    end_line must be the last line of this file - the builder computes it
%    dynamically since 2026-10-14); cards without patchenv never load it.
\renewenvironment{reference}{}{}
\renewenvironment{slogan}{}{}
\renewenvironment{history}{}{}
% empty out citations inside navigation-pointer environments (no bibliography
% backend; spans themselves are cite-free - scanned 2026-10-14, lane-s3)
\renewcommand{\cite}[2][]{}
\usepackage{cleveref}
\crefname{originaltheorem}{original theorem}{original theorems}
\Crefname{originaltheorem}{Original theorem}{Original theorems}
\makeatletter
\newlabel{lemma-diagonal-contractions}{{0GHF}{source}{}{sourcepointer-1}{}}
\newlabel{lemma-diagonal-contractions@cref}{{[originaltheorem][0][]0GHF}{[1][0][]source}{}{}{}}
\makeatother
\crefname{originaltheorem}{original theorem}{original theorems}
\Crefname{originaltheorem}{Original theorem}{Original theorems}
\makeatletter
\newlabel{lemma-finite-dim}{{0GI7}{source}{}{sourcepointer-2}{}}
\newlabel{lemma-finite-dim@cref}{{[originaltheorem][0][]0GI7}{[1][0][]source}{}{}{}}
\makeatother
\crefname{originaltheorem}{original theorem}{original theorems}
\Crefname{originaltheorem}{Original theorem}{Original theorems}
\makeatletter
\newlabel{lemma-formal-object-effective}{{0GI8}{source}{}{sourcepointer-3}{}}
\newlabel{lemma-formal-object-effective@cref}{{[originaltheorem][0][]0GI8}{[1][0][]source}{}{}{}}
\makeatother
\crefname{originaltheorem}{original theorem}{original theorems}
\Crefname{originaltheorem}{Original theorem}{Original theorems}
\makeatletter
\newlabel{lemma-functor-is-solution}{{0GHC}{source}{}{sourcepointer-4}{}}
\newlabel{lemma-functor-is-solution@cref}{{[originaltheorem][0][]0GHC}{[1][0][]source}{}{}{}}
\makeatother
\crefname{originaltheorem}{original theorem}{original theorems}
\Crefname{originaltheorem}{Original theorem}{Original theorems}
\makeatletter
\newlabel{lemma-limit-preserving}{{0GI5}{source}{}{sourcepointer-5}{}}
\newlabel{lemma-limit-preserving@cref}{{[originaltheorem][0][]0GI5}{[1][0][]source}{}{}{}}
\makeatother
\crefname{originaltheorem}{original theorem}{original theorems}
\Crefname{originaltheorem}{Original theorem}{Original theorems}
\makeatletter
\newlabel{lemma-openness-versality}{{0GIA}{source}{}{sourcepointer-6}{}}
\newlabel{lemma-openness-versality@cref}{{[originaltheorem][0][]0GIA}{[1][0][]source}{}{}{}}
\makeatother
\crefname{originaltheorem}{original theorem}{original theorems}
\Crefname{originaltheorem}{Original theorem}{Original theorems}
\makeatletter
\newlabel{lemma-rs}{{0GI6}{source}{}{sourcepointer-7}{}}
\newlabel{lemma-rs@cref}{{[originaltheorem][0][]0GI6}{[1][0][]source}{}{}{}}
\makeatother
\crefname{originaltheorem}{original theorem}{original theorems}
\Crefname{originaltheorem}{Original theorem}{Original theorems}
\makeatletter
\newlabel{lemma-sheaf}{{04FH}{source}{}{sourcepointer-8}{}}
\newlabel{lemma-sheaf@cref}{{[originaltheorem][0][]04FH}{[1][0][]source}{}{}{}}
\makeatother
\crefname{originaltheorem}{original theorem}{original theorems}
\Crefname{originaltheorem}{Original theorem}{Original theorems}
\makeatletter
\newlabel{proposition-spaces-diagonal-representable-noetherian}{{0GE7}{source}{}{sourcepointer-9}{}}
\newlabel{proposition-spaces-diagonal-representable-noetherian@cref}{{[originaltheorem][0][]0GE7}{[1][0][]source}{}{}{}}
\makeatother
\crefname{originaltheorem}{original theorem}{original theorems}
\Crefname{originaltheorem}{Original theorem}{Original theorems}
\makeatletter
\newlabel{remark-separated-needed}{{0GIC}{source}{}{sourcepointer-10}{}}
\newlabel{remark-separated-needed@cref}{{[originaltheorem][0][]0GIC}{[1][0][]source}{}{}{}}
\makeatother
\crefname{originaltheorem}{original theorem}{original theorems}
\Crefname{originaltheorem}{Original theorem}{Original theorems}
\makeatletter
\newlabel{restricted-definition-formal-modification}{{0GDK}{source}{}{sourcepointer-11}{}}
\newlabel{restricted-definition-formal-modification@cref}{{[originaltheorem][0][]0GDK}{[1][0][]source}{}{}{}}
\makeatother
\title{Artin's theorem on contractions}
\author{The Stacks Project Authors}
\date{stacks-project:artin.tex:local-snapshot-20261012; GNU FDL 1.2 or later}
\begin{document}
\maketitle
\noindent\textit{Editorial note.} Statement span is lines 6502-6530 of sources/stacks-artin/pinned-original-source.tex; proof span is lines 6532-6545, the complete author proof (12 body lines).\par\medskip
\noindent\textit{Editorial note.} Navigation stubs: 11 cited-but-not-included labels are supplied as non-load-bearing pointers with their pinned tags where the snapshot tag index has them.\par\medskip
\noindent\textit{Editorial note.} extra\_packages is empty on purpose: the retained preamble already loads xy with options and a second declaration clashes in the assembler wrapper.\par\medskip
\noindent\textit{Editorial note.} Text is retained verbatim, including the author's own typos and placeholders.\par\medskip
\noindent\textit{Editorial note.} Queue index 41 of queue\_artin.json (pb=672, stmt 6502-6530, proof 6532-6545).\par\medskip
\noindent\textit{Editorial note.} stacks-artin wave ar01 (2026-10-14): fresh adjudication accept per lane-s3/artin\_judgments.md; pool per Chief 10:5x direct dispatch (idle-first over S12's 01:5x note, SMC precedent); artin.tex 99 rows/67 with proof judged 46 accept / 21 reject / 5 AGP-overlap avoidance rows; patch-sd-envs contract spans scanned.\par\medskip
\noindent\textit{Editorial note.} Statement carries a reference/slogan/history environment (patchenv -> patch-sd-envs.tex).\par\medskip
\noindent\textit{Allowed standard prerequisites (editorial).}\begin{itemize}
\item Explicit prerequisite boundary. The cited inputs are prerequisites; the argument in the retained span is not replaced by any cited result.
\end{itemize}\medskip
\ifdefined\hypertarget\hypertarget{sourcepointer-1}{}\fi
\noindent\textit{Source reference.} Original theorem 0GHF: sources/stacks-artin/pinned-original-source.tex, source lines 5613--5617. Cited by the retained statement or proof but not part of the collected spans; supplied only as a navigation pointer. The argument does not depend on its content.\par\medskip
\ifdefined\hypertarget\hypertarget{sourcepointer-2}{}\fi
\noindent\textit{Source reference.} Original theorem 0GI7: sources/stacks-artin/pinned-original-source.tex, source lines 6212--6215. Cited by the retained statement or proof but not part of the collected spans; supplied only as a navigation pointer. The argument does not depend on its content.\par\medskip
\ifdefined\hypertarget\hypertarget{sourcepointer-3}{}\fi
\noindent\textit{Source reference.} Original theorem 0GI8: sources/stacks-artin/pinned-original-source.tex, source lines 6237--6240. Cited by the retained statement or proof but not part of the collected spans; supplied only as a navigation pointer. The argument does not depend on its content.\par\medskip
\ifdefined\hypertarget\hypertarget{sourcepointer-4}{}\fi
\noindent\textit{Source reference.} Original theorem 0GHC: sources/stacks-artin/pinned-original-source.tex, source lines 5372--5377. Cited by the retained statement or proof but not part of the collected spans; supplied only as a navigation pointer. The argument does not depend on its content.\par\medskip
\ifdefined\hypertarget\hypertarget{sourcepointer-5}{}\fi
\noindent\textit{Source reference.} Original theorem 0GI5: sources/stacks-artin/pinned-original-source.tex, source lines 5636--5640. Cited by the retained statement or proof but not part of the collected spans; supplied only as a navigation pointer. The argument does not depend on its content.\par\medskip
\ifdefined\hypertarget\hypertarget{sourcepointer-6}{}\fi
\noindent\textit{Source reference.} Original theorem 0GIA: sources/stacks-artin/pinned-original-source.tex, source lines 6452--6455. Cited by the retained statement or proof but not part of the collected spans; supplied only as a navigation pointer. The argument does not depend on its content.\par\medskip
\ifdefined\hypertarget\hypertarget{sourcepointer-7}{}\fi
\noindent\textit{Source reference.} Original theorem 0GI6: sources/stacks-artin/pinned-original-source.tex, source lines 6175--6178. Cited by the retained statement or proof but not part of the collected spans; supplied only as a navigation pointer. The argument does not depend on its content.\par\medskip
\ifdefined\hypertarget\hypertarget{sourcepointer-8}{}\fi
\noindent\textit{Source reference.} Original theorem 04FH: sources/stacks-artin/pinned-original-source.tex, source lines 5625--5629. Cited by the retained statement or proof but not part of the collected spans; supplied only as a navigation pointer. The argument does not depend on its content.\par\medskip
\ifdefined\hypertarget\hypertarget{sourcepointer-9}{}\fi
\noindent\textit{Source reference.} Original theorem 0GE7: sources/stacks-artin/pinned-original-source.tex, source lines 5072--5083. Cited by the retained statement or proof but not part of the collected spans; supplied only as a navigation pointer. The argument does not depend on its content.\par\medskip
\ifdefined\hypertarget\hypertarget{sourcepointer-10}{}\fi
\noindent\textit{Source reference.} Original theorem 0GIC: sources/stacks-artin/pinned-original-source.tex, source lines 6548--6569. Cited by the retained statement or proof but not part of the collected spans; supplied only as a navigation pointer. The argument does not depend on its content.\par\medskip
\ifdefined\hypertarget\hypertarget{sourcepointer-11}{}\fi
\noindent\textit{Source reference.} Original theorem 0GDK: sources/stacks-artin/pinned-original-source.tex, source lines 6545--6545. Cited by the retained statement or proof but not part of the collected spans; supplied only as a navigation pointer. The argument does not depend on its content.\par\medskip
\begin{theorem}
\label{theorem-contractions}
\begin{reference}
\cite[Theorem 3.1]{ArtinII}
\end{reference}
Let $S$ be a locally Noetherian scheme such that $\mathcal{O}_{S, s}$
is a G-ring for all finite type points $s \in S$. Let $X'$ be an algebraic
space locally of finite type over $S$. Let $T' \subset |X'|$ be a closed
subset. Let $W$ be a locally Noetherian formal algebraic space over $S$
with $W_{red}$ locally of finite type over $S$. Finally, we let
$$
g : X'_{/T'} \longrightarrow W
$$
be a formal modification, see Algebraization of Formal Spaces, Definition
\ref{restricted-definition-formal-modification}. If $X'$ and $W$ are
separated\footnote{See Remark \ref{remark-separated-needed}.} over $S$, then
there exists a proper morphism $f : X' \to X$ of algebraic spaces over $S$,
a closed subset $T \subset |X|$, and an isomorphism $a : X_{/T} \to W$
of formal algebraic spaces such that
\begin{enumerate}
\item $T'$ is the inverse image of $T$ by $|f| : |X'| \to |X|$,
\item $f : X' \to X$ maps $X' \setminus T'$ isomorphically to
$X \setminus T$, and
\item $g = a \circ f_{/T}$ where $f_{/T} : X'_{/T'} \to X_{/T}$
is the induced morphism.
\end{enumerate}
In other words, $(f : X' \to X, T, a)$ is a solution as defined earlier in
this section.
\end{theorem}

\begin{proof}
Let $F$ be the functor constructed using $X'$, $T'$, $W$, $g$ in this section.
By Lemma \ref{lemma-functor-is-solution} it suffices to show that
$F$ corresponds to an algebraic space $X$ locally of finite type over $S$.
In order to do this, we will apply
Proposition \ref{proposition-spaces-diagonal-representable-noetherian}.
Namely, by Lemma \ref{lemma-diagonal-contractions}
the diagonal of $F$ is representable by closed immersions
and by
Lemmas \ref{lemma-sheaf}, \ref{lemma-limit-preserving},
\ref{lemma-rs}, \ref{lemma-finite-dim},
\ref{lemma-formal-object-effective}, and \ref{lemma-openness-versality}
we have axioms [0], [1], [2], [3], [4], and [5].
\end{proof}

\end{document}
